% !TeX root = ../../../main.tex
\subsubsection{Forecasting Tasks}
Using the timelines extracted from the annotations, we defined two primary forecasting tasks and a survival analysis task as follows:

\ChSevenParagraph{Event Occurrence Prediction.} 
Given a prefix of the clinical timeline (all events up to a certain time $t$), the model is tasked with predicting whether each of the immediate next $k$ events occurs within a specified time horizon. 
This setup is repeated across multiple time cut-offs to simulate an “online” forecasting scenario, where the model must output a binary label for each of the next $k$ events:
does this event occur within $h$ hours after $t$? 
Time horizons used include 1 hour, 24 hours (1 day), and 168 hours (1 week). 
The task is framed as a series of binary classification problems, with one binary decision per event. 
Evaluation is based on precision, recall, and F1 score, averaged across event positions for each time horizon.

\ChSevenParagraph{Temporal Ordering Prediction.} 
This task assesses the model's ability to reconstruct the correct sequence of future events. 
At each time cut-off $t$, we extract the next $k$ events from the timeline and remove their timestamps. 
The model must output a permutation of these events that matches their true chronological order.
This is framed as a ranking task, evaluated by computing the pairwise concordance between the predicted and true orderings (e.g., proportion of correctly ordered pairs). 
This tests whether the model can infer temporal progression from unordered event content alone.

\ChSevenParagraph{Survival Analysis for Mortality Time.}  
We include a classical survival analysis task to model time until death. 
Many case reports specify whether the patient died and, if so, when (e.g., "the patient died on hospital day 10"), which enables us to define a time-to-event outcome. 
For this task, we evaluate models at predefined cut-off times--- specifically at 0 hours (admission), 24 hours (1 day), and 168 hours (1 week)--- and use the extracted event history up to each cut-off as input. 
A survival model is trained to predict the probability of survival over time beyond each cut-off. 
We evaluate using the time-dependent concordance to measure alignment of the predicted survival times with actual outcomes.

